% FILE: 06_contribution_to_thesis.tex
% Project: standards
% Thesis Role: public-standards-compliance-and-gravity-field-mapping

\section{Contribution to the Dissertation}
\label{sec:standards-contribution}

\subsection{Contribution to Prediction vs.~Coordination}
\label{sec:standards-contribution-prediction}

The dissertation's central theoretical claim is that the transition from prediction to
coordination requires shared formal referents---structures that two independent parties
can both evaluate and agree upon without trusting each other's good faith. The standards
project is the primary evidence that this transition is achievable in the domain of
process intelligence.

By mapping 16 public standards to named mathematical formulations and named validation
thresholds, the project demonstrates that process conformance claims can be made
coordination-capable: a buyer's advisor and a seller's auditor can both execute the same
conformance query on the same XES log against the same WF-net model and arrive at the
same fitness score within a $10^{-6}$ tolerance. This is coordination. Neither party
needs to trust the other's assertion; each can independently evaluate the shared referent.

The Reverse Porter's Five Forces analysis (\texttt{reverse-porter-five.md}) articulates
the competitive consequence of this shift: once process evidence is coordination-capable,
the market structure reorganizes around the standard keeper. The five forces are not
merely blunted---they are inverted. Threat of new entrants becomes weaponized entry
channeling; rivalry becomes consensual conformance. This reorganization is only possible
because the standards project has established the shared referent that makes coordination
possible in the first place.

\subsection{Contribution to Process Evidence}
\label{sec:standards-contribution-evidence}

The standards project contributes the \emph{external anchoring} of the Evidence lattice
\texttt{Evidence<T, State, Witness>}. The lattice is a type-system construct; without
external anchoring, it witnesses only internal consistency. The 16 standard placements
provide the anchors: each witness marker in \texttt{src/witness.rs} of the wasm4pm-compat
crate corresponds to a named public standard, and each standard's placement file defines
what the witness certifies.

Concretely, the OCEL 2.0 placement defines that the \texttt{Ocel20} witness certifies
bipartite graph validity: $\forall (e, o, q) \in \text{E2O},\ e \in \mathcal{E} \land
o \in \mathcal{O}$. The WF-net soundness placement defines that the
\texttt{WfNetSoundnessPaper} witness certifies option-to-complete, proper completion, and
no dead transitions per van der Aalst (1998). The \texttt{Xes1849} witness certifies
trace hash chaining and chronological invariance. These are compile-time invariants in
Rust, not runtime assertions---the type system enforces that evidence claiming a standard
carries the correct witness type. The gravity field is not merely a marketing claim; it
is a compile-time invariant, as stated in \texttt{PROCESS\_STANDARDS\_GRAVITY\_FIELD.md}.

\subsection{Contribution to Manufacturing Architecture}
\label{sec:standards-contribution-manufacturing}

The standards project contributes two architectural layers to the manufacturing pipeline.

\paragraph{The ggen projection map.}
The file \texttt{public\_standards\_to\_ggen\_projections.md} defines the mapping from
standard constructs to Rust code artifacts materialized by ggen: Petri net places project
to Rust structs, BPMN gateways project to \texttt{match}/\texttt{join!} patterns, and
Declare templates project to FSA implementations. This map is the bridge from the
standards registry to the wasm4pm execution surface. Without it, ggen would be a code
emitter without a formal basis; with it, every emitted artifact is traceable to a named
standard.

\paragraph{The MAPE-K integration.}
The MAPE-K integration file (\texttt{MAPE\_K\_INTEGRATION.md}) defines how the standards
architecture maps to the four phases of the autonomic feedback loop. Monitor maps to
streaming OTel trace ingestion; Analyze maps to streaming conformance checking via
\texttt{StreamingConformanceSession}; Plan maps to prescriptive process intelligence via
\texttt{prediction.rs} outcome models and \texttt{diagnostic.rs} remediation shapes;
Execute maps to \texttt{OcelEvent} actuation. This mapping makes the manufacturing
architecture self-monitoring: the same standards that define valid process outputs also
define the monitoring signals that detect deviations from those outputs.

\subsection{Contribution to Capital-Flow Infrastructure}
\label{sec:standards-contribution-capital}

The standards project contributes the legal and cryptographic basis for M\&A-admissible
process evidence. The Slide-to-Receipt bridge formula
$\text{SlideBinding} = \operatorname{BLAKE3}(\text{SlideUUID} \parallel \text{ReceiptHash}
\parallel \text{ClaimValue})$ makes every slide in an investment presentation carrying an
operational claim independently verifiable by the buyer's advisors. This is not a
documentation convenience; it is a structural requirement for board admissibility as
defined in \texttt{ma/define\_board-admissible\_claim\_requirements.md}.

The reverse lock-in doctrine extends this contribution to the capital-flow level. By
establishing $W_{\text{std}}$ as a publicly archived, cryptographically signed, sound WF-net,
the standards project creates the conditions for governance rent extraction: standard
updates, proof generation, proof validation, conformance certification, and market
intelligence are all services that rational agents in the verified economy must purchase
from the standard keeper. The Nash equilibrium cage analysis in \texttt{reverse-lock-in.md}
formalizes this as a game-theoretic outcome: rational agents choose conformance because
non-conformance is mathematically unprofitable, and the standard becomes the operating
system of the verified process economy.
